% 第1章第7节：本章小结。

\section{本章小结}
\label{sec:fs-summary}

章首的问题是：当表示规模不断变化时，怎样直接比较函数、谈论极限和正交分量。本章的回答可以归结为四个依次进行的选择：选定函数集合，确定什么算作同一个元素，规定误差，检查极限是否留在集合中。

第一个选择决定线性运算是否封闭：连续函数、多项式和固定函数的张成是线性空间，参数化的非线性族一般不是；一维输入也可以产生无限维的函数空间。第二个选择决定元素的身份：用积分衡量误差时，只在零测集上不同的函数必须看成同一个元素，代价是等价类一般没有确定的点值。第三个选择是范数，它让误差满足正定、齐次和三角不等式，从而接入度量空间的全部工具。第四个选择是完备性检查，它同时依赖集合和范数：同一个$C([-1,1])$在一致范数下完备，在平方积分范数下不完备，斜坡逼近阶跃的例子说明了二者的差别，而$L^2([-1,1])$正是后者的完备化。

三类空间的结构可以这样对照。连续函数空间$C(K)$的元素是普通函数，一致范数控制每一点，完备性来自一致收敛保持连续。$L^p(\mu)$的元素是等价类，范数累计整体偏差，Hölder与Minkowski不等式保证范数公理，完备性靠抽子列与Fatou引理（有限$p$）或合并零测例外集（$p=\infty$）得到。数列空间$\ell^p$是计数测度下的$L^p$，等价类退化为数列本身。比较这些空间时，要分清三种保证：同一个函数集合不保证相同的收敛；相同的零集保证换测度后元素的识别一致，却不保证范数可比；只有与序号无关的范数不等式才能传递收敛。三个对照序列$x^n$、窄尖峰和打字机序列标出了一致、$L^p$与逐点收敛之间不能推出的方向，子列定理与控制收敛则给出了在附加条件下可以推出的部分。

Hilbert空间在完备性之外还有内积。闭子空间上的正交投影存在且唯一，由残差正交刻画，并给出分解$H=M\oplus M^\perp$；闭性缺失时最近点可能不存在。正交归一系给出坐标：Bessel不等式对任何正交系成立，Parseval等式只对完备正交系成立，截断误差等于系数平方的尾和。但$L^2$意义下的展开收敛并不推出逐点或一致收敛。

本章把函数当作空间中的点。接下来的问题是作用在这些点上的映射：哪些线性映射能稳定地变换函数，标量读取能否统一地表示，投影作为映射有哪些性质，以及当整体误差不趋于零时，能否通过所有连续线性读取来判断极限。这些问题由第~\ref{chap:linear-functionals-operators}章的线性算子、线性泛函、对偶与弱收敛处理。
